\section{From a missing chapter to a missing possibility}
\label{sec:opening}
A reader asks a frontier language model to summarize a book. The summary is fluent, balanced in tone, and apparently complete. The reader has not read the book, so the very quality that makes the summary convenient also makes it difficult to verify. A second model receives the book and the summary and is asked to identify a consequential omission. It points to a particular passage, explains why that passage changes the conclusion, and provides enough context for an independent check.

The assessment problem has changed. The reader can investigate a specific discrepancy instead of trying to certify the entire compression from its style. This asymmetry is useful: one decisive example can refute a particular claim even when constructing or evaluating the whole artifact is difficult. The second model can nevertheless be wrong. It may invent a passage, misread a qualification, or complain that a short summary does not contain everything. The task should therefore ask for the strongest \emph{supported} criticism and allow the answer that no consequential defect was established.

This observation has direct precedents. Model-written critiques helped human evaluators find flaws in topic-based summaries in the experiments of Saunders et al.\ \citep{saunders}. Later work on code critics found both useful bug detection and hallucinated criticisms \citep{critic}. The opening of \emph{Transformational Intelligence} already develops the book example, adds independent adjudication, and considers producing a source-grounded structure before compression \citep[\S1 and Appendix C]{boucher_ti}. We build on that foundation rather than claim the second critic as a new invention.

\subsection{The larger omission}
An enterprise faces a related problem with a crucial difference. A book has a finite source. The world does not supply a complete volume containing every possible invention, production process, or voluntary arrangement. A proposal can accurately describe its candidate and still fail to consider a much better route. A scientific investigation can compare its hypotheses carefully while missing an instrument that would change which questions are answerable. An economically sensible rejection can omit a complementary construction that would make the activity viable.

The central research question is therefore:
\begin{thesisbox}
\textbf{Can intelligence turn a consequential, verified gap in what we know or consider into the means of doing something valuable that the current arrangement cannot deliver?}
\end{thesisbox}
The answer requires more than another round of disagreement. A valid objection has to connect to a diagnosis, a response, an executable construction, and a complete comparison. Some objections call for correction. Others call for new evidence. A third class reveals a removable dependency: the current method requires something expensive or unavailable, but an alternative may achieve the same authorized outcome without it. The best response can also be procurement, an established method, or no further action.

Consider a hypothetical small manufacturer that rejects a new product because qualifying each short production run appears too expensive. A failure challenger asks whether its proposed shortcut really establishes the required quality. A possibility challenger asks whether several runs share a testable structure that would make a common qualification instrument worthwhile. An instrument challenger asks whether the existing measurements can distinguish acceptable from unacceptable behavior. These are different questions. Their answers might support a shared test capability, a different production method, or the original rejection. The useful advance is whatever survives construction and evidence.

\subsection{How the three papers fit together}
\begin{table}[htbp]
\small
\begin{tabularx}{\textwidth}{P{1.48in}YY}
\toprule
\textbf{Foundation} & \textbf{What it makes editable} & \textbf{Its role here}\\
\midrule
Transformational Intelligence & Representations, information, dependencies and intermediate means, with the authorized purpose preserved. & Diagnose whether a criticism exposes an avoidable restriction; construct and test another formulation.\\
Alpha Genesis Flywheel $\Omega$ & The conversion of frontier capability into useful possibilities and better discovery methods. & Replay consequential gaps when a relevant capability changes; distinguish better tools from better invention.\\
Invent the Economy & Productive capability and the arrangement of participants, transfers, risk and timed funding. & Build the complete conditions under which useful work is technically valid and voluntarily worthwhile.\\
\bottomrule
\end{tabularx}
\caption{An explicit synthesis of the supplied manuscripts \citep{boucher_ti,boucher_genesis,boucher_economy}. The organizing transition of this manuscript is the move from an adjudicated challenge to a tested productive response.}
\end{table}

\subsection{Contribution and intellectual position}
The paper contributes an operational \emph{Missing Move contract}; a distinction among challenges to an answer, its opportunity set, and its evidence instrument; conditional accounting for useful and harmful criticism; and a matched evaluation programme connecting constructive challenge to frontier adaptation and economic continuation. An executable finite case carries the central transition through: find a consequential information loss, construct a sufficient representation, verify its minimality, and test a changed successor. A separate cash-path counterexample shows why a profitable successful delivery can still leave a promised refund unfunded. The larger architecture remains a proposal.

Counterexample-guided synthesis provides an established model for generating candidates, finding failures, and using those failures to constrain subsequent construction \citep{cegis}. Debate offers a different route to eliciting useful information for a judge \citep{debate}. Neither formal finite-program guarantees nor idealized debate arguments make open-world invention automatically reliable. A natural-language challenge may lack a sound specification, a complete checker, or even an adequate measurement. We make those missing conditions part of the work.

The ambition is expansive without requiring every component to be novel. An ordinary test, a familiar algorithm, or a conventional contracting arrangement can be the decisive missing move. Scientific novelty, useful construction, commercial success, and improvement of discovery have different evidence requirements. A competent institution pursues the valuable result and states precisely which of those achievements it has established.

\section{The Missing Move contract}
\label{sec:contract}
\subsection{Keep purpose, means and measurement distinct}
Let the purpose contract be
\begin{equation}
\I=(O,H,A,B,T),
\end{equation}
where $O$ denotes the intended outcome, $H$ the hard requirements, $A$ the authority to act, $B$ the available resources, and $T$ the relevant time conditions. A formulation $F$ describes a particular set of representations, information procedures, allowed internal actions and intermediate targets. An evaluation instrument $E$ supplies evidence about whether a realization serves $\I$. This separation follows the foundational manuscript \citep[\S2]{boucher_ti}.

A critic may propose changing $F$ or improving $E$. It cannot silently amend $O$, weaken $H$, expand $A$, or spend imaginary resources. If the required deliverable is a two-page summary, an interactive application does not fulfill the instruction merely because it might be more interesting. If the purpose is to understand an argument before a meeting, the application may be a legitimate alternative. The distinction depends on the actual assignment. A broad discovery charter may authorize proposing new beneficiary outcomes within a stated mission. Once a candidate enters confirmation, freeze its concrete outcome contract; a material change of ends or permissions begins an explicitly amended comparison.

\begin{definition}[Consequential challenge]
Relative to a stated purpose and decision, a consequential challenge is an allegation that a candidate claim, omitted condition, excluded alternative, or inadequate evidence procedure could materially change the required outcome or a justified choice of action. It is a hypothesis until its relevant evidence and materiality are adjudicated.
\end{definition}
An omission is therefore not automatically a defect. A summary must compress; an engineering plan must abstract; a research portfolio must leave possibilities unexplored. Materiality depends on the user's purpose, the decision at stake, and the available budget. Exhaustive commentary can make a useful artifact worse.

\begin{seabox}
\textbf{A missing distinction you can inspect.} A report stores only the mean of two readings and claims it can determine whether every reading is at most 90. The traces $(60,100)$ and $(80,80)$ have the same mean, 80, but opposite acceptance outcomes. That pair is a concrete witness against the claim. Retaining the maximum before compression is an ordinary sufficient response; it does not recover an already discarded maximum. Section~\ref{sec:construction} makes the transition from this witness to a constructed representation and a changed successor executable. The point is visible: a missing distinction becomes a specific construction target.
\end{seabox}

\subsection{Two linked records: the challenge and the construction}
We separate a challenge record $w$ from a response record $r$. This prevents the correctness of a diagnosis from being mistaken for the correctness of its favorite remedy:
\begin{align}
 w &= (c,s,e,m,q),\\
 r &= (\rho,F',k,\tau,v).
\end{align}
Here $c$ identifies the challenged claim or choice; $s$ fixes its scope and source versions; $e$ is inspectable evidence; $m$ states decision materiality; and $q$ gives a falsifying or discriminating check. The response identifies the diagnosed restriction $\rho$, alternative formulation $F'$, executable construction $k$, complete test $\tau$, and validity and invalidation conditions $v$.

For a book, $e$ may be an accurately located passage with enough surrounding context to preserve its meaning. For a program, it can be an executable input and failing output. For a market arrangement, it may be a feasible participant deviation, an omitted liability, or a cash-flow prefix that cannot be funded. For a scientific instrument, the evidence may show that two materially different states are observationally indistinguishable under the proposed measurement procedure. Each domain needs its own semantics of evidence.

\begin{table}[htbp]
\small
\begin{tabularx}{\textwidth}{P{1.25in}YY}
\toprule
\textbf{Question} & \textbf{Minimum useful answer} & \textbf{What it cannot establish alone}\\
\midrule
What is challenged? & An exact assertion, required property or excluded route. & That the complete candidate is bad.\\
Where is the evidence? & Source identity and location, a runnable witness, or a justified measurement. & That the evidence source is infallible.\\
Why does it matter? & A specified effect on purpose, acceptance or decision. & That every omitted detail deserves repair.\\
What kind of obstacle? & A hard constraint to honor, an uncertain premise to resolve, or a restriction of means to change. & Permission to weaken a real requirement.\\
What is the response? & A buildable alternative plus a test against serious rivals. & That a valid criticism makes the repair valid.\\
When does it transfer? & Preconditions, costs, interfaces and invalidation rules. & Current acceptance or authority for a successor.\\
\bottomrule
\end{tabularx}
\caption{A compact contract for constructive challenge. Evidence and construction are linked while remaining separately accountable.}
\end{table}

\subsection{Three directions of challenge}
\textbf{The failure challenger} searches for a supported reason the proposed route would fail: a contradiction, missing exception, unacceptable operating case, uncompensated participant, or unpriced obligation. Its best output is a decisive witness, not a large number of objections.

\textbf{The possibility challenger} constructs a serious alternative the current candidate set excludes. This matters because a selector can choose perfectly among bad proposals. The relevant comparison is between complete routes to the same intended outcome, including an existing solution or doing nothing. A purportedly impossible activity becomes a candidate opportunity only when a concrete response changes the constraint responsible for that judgment.

\textbf{The instrument challenger} examines whether the available evidence can settle the dispute. It may propose a new observation, an independent implementation, a different decomposition, or a validated measure closer to the intended outcome. Its product is a better way to find out. It must itself be tested, with both sides rescored under a common supported instrument when that instrument changes \citep[\S7]{boucher_ti}.

These are functions, not necessarily three agents. One capable system, several systems, a domain expert or conventional software may perform them. Role separation helps organize work; it does not manufacture independent errors or independent evidence.

\subsection{Adjudicate before repair, and compare several responses}
An adjudicator returns one of four useful states: \emph{supported and material}, \emph{supported but immaterial}, \emph{refuted}, or \emph{unresolved}. ``Unresolved'' is not a softer form of acceptance. It identifies an uncertainty that the next decision must either resolve, tolerate explicitly, eliminate through a changed construction, or avoid.

A supported material challenge opens a response set:
\begin{equation}
\mathcal R(w)=\{\text{repair},\ \text{measure},\ \text{construct},\ \text{rearrange},\ \text{adopt},\ \text{defer or stop}\}.
\end{equation}
Assess critics by independently verified material findings and downstream effects, including harmful revisions and wasted effort. Rewarding objection count, rhetorical force or frequency of overturning the producer encourages the wrong behavior. The categories may overlap. A shared instrument and a staged commitment can be jointly useful even when neither alone justifies its cost. A system that immediately repairs the first valid objection can miss a simpler route that removes the disputed dependency altogether. Conversely, adding an elaborate institution to solve a local defect can be worse than a small correction.

\begin{seabox}
\textbf{The constructive question is not only ``What is wrong?''}\par
It is: \textbf{``What would have to become true, available, measurable or mutually worthwhile for the valuable outcome to work---and can we make that happen?''}
\end{seabox}

\subsection{An opportunity witness needs two sides}
A counterexample can demonstrate that a universal assertion is false on a stated case. An opportunity claim needs more. It needs a \emph{positive side}: a complete route that meets the purpose under its stated conditions. It also needs a \emph{comparative side}: why this route matters relative to available alternatives, including a route that reconstructs the supposed missing asset conventionally.

Let $J_\I(x)$ be a declared net-outcome measure for an admissible complete route $x$, and let $\B$ be a disclosed finite, nonempty comparison set. Define
\begin{equation}
\alpha_\I(x;\B)=J_\I(x)-\max_{b\in\B}J_\I(b).
\end{equation}
This is a comparison to $\B$, not a proof about every possible enterprise. If outcomes cannot legitimately be collapsed into one scalar, retain a vector of relevant outcomes and report the actual trade-offs. Hard constraints remain constraints rather than purchasable score penalties.

The strongest formulation-dependent claim adds a justified exclusion: the old setup cannot attain the specified outcome within its defined permissions, information and resources. In real research that exclusion is often unavailable. Then the honest claim is a supported gain over the strongest tested complete alternatives. Neither an unsuccessful search nor a pessimistic first draft proves impossibility.

\section{When is another critic worth its cost?}
\label{sec:value}
\subsection{A conditional accounting identity}
Fix one pre-specified, decision-material defect proposition per episode, relative to a stated purpose and task distribution. Let $F$ denote that this exact proposition is true. Let $D$ denote a criticism alleging this same defect, and let $A_c\subseteq D$ denote acceptance of that allegation followed by its specified response. A flawed candidate can receive a false criticism about a different issue; such mixed-target episodes require a richer event partition and are outside this binary reference model. Define
\begin{align}
\pi&=\Prob(F), & d&=\Prob(D\mid F), & f&=\Prob(D\mid\neg F),\\
 a&=\Prob(A_c\mid D,F), & b&=\Prob(A_c\mid D,\neg F).
\end{align}
The letter $A_c$ here refers to an accepted challenge, not the authority component of $\I$. Let $G=\E[\Delta J\mid F,D,A_c]$ be the expected net utility improvement conditional on accepting a true allegation of this target defect; let $H=-\E[\Delta J\mid\neg F,D,A_c]$ be the expected loss conditional on accepting a false allegation. Both include revision and retesting costs on their respective branches. Let $K$ be the expected incremental cost of criticism and adjudication across all branches, including retrieval, discarded objections, latency and opportunity cost. On a zero-probability conditioning event, use a finite declared kernel value: its weighted contribution is zero, and the distribution does not identify it.

\begin{proposition}[Value of adjudicated challenge]
\label{prop:value}
Suppose branches without an accepted challenge retain incumbent utility before review costs. With the conditional quantities above,
\begin{equation}
\boxed{\Delta U=\pi daG-(1-\pi)fbH-K.}
\label{eq:value}
\end{equation}
For $G>0$ and $H,K\geq0$, the challenge route improves expected utility exactly when
\begin{equation}
\pi daG>(1-\pi)fbH+K.
\label{eq:break}
\end{equation}
\end{proposition}
\begin{proof}
Partition accepted revisions into true and false criticisms. Their probabilities are $\pi da$ and $(1-\pi)fb$ by the chain rule. Multiply by their conditional net effects and subtract the expected review cost. No statistical independence between the producer, critic and adjudicator is required.
\end{proof}

Equation~\eqref{eq:value} is an accounting identity under explicit assumptions. It does not tell us that the inputs are calibrated. Multiplying two models' published marginal benchmark accuracies would not estimate these conditional events. In applications, one must measure the actual joint process on the relevant distribution. A correct diagnosis may lead to a harmful repair, in which case $G$ can be negative and the same accounting identity still applies.

The equation exposes several productive choices. Better source access may improve $d$. A clearer witness may improve adjudication. A smaller repair may reduce harm or cost. A different formulation may reduce the need for recurring criticism. A more expensive critic is justified only when its complete consequences improve the choice. The model therefore supports testing evidence-preserving production and conventional inspection as well as adding models.

\subsection{Confirmation is not coverage}
The probability of a material flaw escaping an \emph{accepted challenge} is
\begin{equation}
\Prob(F\cap\neg A_c)=\pi(1-da).
\label{eq:escape}
\end{equation}
This is not the post-repair defect rate: repairs may fail, and false criticisms may introduce new defects. It isolates a different weakness. Even a perfect adjudicator cannot confirm a flaw that the challenger never proposes. High precision among accepted criticisms can coexist with poor coverage of consequential omissions.

For $m$ critics, detection is $\Prob(\bigcup_{i=1}^{m}D_i\mid F)$. Without a conditional independence assumption, only the basic bounds
\begin{equation}
\max_i\Prob(D_i\mid F)\ \leq\ \Prob\!\left(\bigcup_iD_i\mid F\right)\ \leq\ \min\!\left\{1,\sum_i\Prob(D_i\mid F)\right\}
\end{equation}
hold. Perfectly correlated critics obtain no detection gain by replication. Different model providers may supply useful diversity, but their names do not establish the relevant joint distribution.

An illustrative two-critic calculation makes this concrete. Suppose each critic misses a flaw with probability $m_0$, and the Bernoulli miss indicators have a feasible Bernoulli correlation $\rho$. Then
\begin{equation}
\Prob(\text{both miss})=m_0^2+\rho m_0(1-m_0).
\label{eq:correlation}
\end{equation}
For the deliberately supplied values $m_0=0.30$, independence gives $0.09$; $\rho=0.80$ gives $0.258$. These are mathematical examples, not measured model rates. A shared source omission or common mistaken assumption can place the operation much closer to the second situation.

\subsection{The verifier can also cost more than it contributes}
Compare adjudication with accepting every criticism. In the reference model, revision effects are fixed within each truth class, so accepting or rejecting a criticism does not select a different conditional repair payoff. Under that additional assumption, the gain from selective acceptance before the verifier's incremental cost is
\begin{equation}
(1-\pi)f(1-b)H-\pi d(1-a)G.
\end{equation}
The first term is harm avoided by rejecting false criticisms. The second is useful repair foregone by rejecting true ones. In a more general population, replace $H$ and $G$ in this comparison by $H_{\mathrm{rej}}$ and $G_{\mathrm{rej}}$, the effects of acting on rejected false and true criticisms. Selective adjudication can change these conditional effects; the accepted-branch quantities alone do not identify them. Subtract additional adjudication cost to evaluate the complete route. Verification is not an automatically positive extra stage; the appropriate degree depends on consequence, evidence and cost. Hard acceptance requirements still apply even when an expected-value shortcut looks attractive.

For repeated challenge, the relevant quantity is the conditional value of the \emph{next} permitted action given the evidence already obtained. Continue when its expected decision value exceeds its complete cost, subject to the purpose and obligations. Stop or adopt another means when no supported available action meets that condition. This is a decision rule under beliefs, not a certificate that no valuable unknown action exists. It prevents an infinite regress of critics criticizing critics from becoming the default product.

\section{Invent the evidence before inventing the verdict}
\label{sec:instrument}
\subsection{The world may be missing from the conversation}
A second model that sees only the first model's summary is poorly placed to expose a missing chapter. It may notice internal inconsistency, but it cannot establish source coverage from absent evidence. Similarly, several agents inspecting the same incomplete engineering model cannot validate a physical property that the model does not represent. The foundational paper's passive-observation example makes this limitation exact within a specified causal model \citep[\S3]{boucher_ti}.

There are four useful responses: obtain a discriminating observation; construct a new observation procedure; choose an acceptable route robust to the supported possibilities; or redesign the work so that the uncertain condition no longer matters. Sometimes stopping is appropriate. All four require comparing their full cost with the value of the resulting decision. A longer exchange of opinions is useful only insofar as it changes an adequately grounded inference or creates a better action.

\subsection{Evidence-preserving production is a strong alternative}
Instead of writing a free-form output and reconstructing its basis afterward, the producer can preserve source identity, supporting passages, dependencies, qualifications and unresolved questions as it works. A summary can be compiled from an argument structure. A design can carry a bill of assumptions and executable tests. A commercial proposal can carry its cash schedule and participant outside options.

This construction changes the verification workload, but introduces its own vulnerability: the evidence map can omit the same crucial issue as the answer. The challenger therefore needs access to the original source or environment, not only the producer's selected excerpts. A separate coverage check asks whether the representation itself excluded something material. Both source acquisition and maintenance of the evidence structure belong in the comparison.

Prover--verifier research shows that checkability can be an explicit optimization target in a studied mathematical-reasoning setting \citep{prover_verifier}. It does not turn natural-language justifications into formal proofs. The appropriate artifact is a compact record that makes the decisive claim easy to test while preserving enough context to resist misleading extraction.

\subsection{Challenge the checker without moving the goalposts}
A checker may be complete inside an inadequate class. A reference answer may be wrong. Several apparently independent judgments may depend on the same source. The related \emph{Out-Reflect} manuscript constructs bounded evidence correction for a specified mathematical family and preserves counterexamples to universal self-certification \citep{boucher_reflect}. Its guarantees do not transfer to arbitrary language or markets.

When a material checker defect is established, version the instrument and state the reason for change. Validate the replacement against independent cases, retain the earlier outcome, and rescore both candidate and comparator under the same supported measure. Do not privately replace an unfavorable judge with a favorable one. A new measurement epoch is compatible with scientific progress; retroactively changing the intended outcome is not evidence of improvement.

Operational independence is concrete. It includes separate implementation paths where useful, raw-source access, hidden confirmation cases, blinded arm identity, different measurement mechanisms, and checks against coherent errors. It does not require pretending that every source is independent. Known dependencies should be recorded, and unresolved dependencies should remain visible.

\begin{thesisbox}
\textbf{A criticism is a hypothesis. A witness makes it checkable. A construction makes it useful. A fresh comparison makes its advantage credible.}
\end{thesisbox}
